\section{AI 使用、试错过程与人工采用}\label{app:ai}
\subsection{使用范围与责任}
AI 辅助研究问题梳理、文献检索与摘录、程序开发、推导检查、候选认知编码、判定说明和图文表达。研究团队确定问题、假设、指标及解释，核对原文、公式和输出，并决定最终采用内容。批量测评由 Next 与 TokenDance 平台提供接口；Next 为赛事 Token 指南所列算力入口。

\begin{table}[htbp]\centering\small
\caption{主要 AI 工具及实际用途}\label{app:ai-tools}
\begin{tabularx}{\textwidth}{@{}p{.18\textwidth}p{.32\textwidth}X@{}}\toprule
环节 & 工具与精确标识 & 用途\\\midrule
研究与实现 & Codex，\path{gpt-6-astra} & 问题和文献梳理、建模与代码辅助、论文与结果核对。\\
基础编码 & DeepSeek V4.1 Flash，\path{deepseek-v4.1-flash} & 独立判断任务与已展示贡献，输出连续原文、解释及弃权。\\
基础编码 & GLM 5.3 Flash，\path{glm-5.3-flash} & 相同输入和规则下的第二独立基础判断。\\
独立复核 & Claude Sonnet 5，\path{claude-sonnet-5} & 对固定随机与高风险子集作独立复核。\\
独立复核 & GPT-6 Sol，\path{gpt-6-sol} & 同一复核子集的第二独立判断。\\
说明性插图 & ChatGPT 图像生成 & 依据团队提供的文字和技术关系生成说明图；团队核查内容。\\\bottomrule
\end{tabularx}
\end{table}
模型身份由实际调用计划和记录核对。基础编码使用 TokenDance，复核使用 Next；早期三模型文本预实验另使用 Gemini 2.5 Pro（\path{gemini-2.5-pro}）。服务端细分版本以接口记录为准。本文定量图由聚合分析代码绘制，图像生成服务用于说明性材料。

\subsection{从候选判断到人工采用的探索链}
探索过程按输入、AI 产物、可见问题、核查与采用记录串联。文献阶段由 AI 提取研究设计、样本和结论，再逐页查回来源；曾在研究资料整合中发现 $8+43+55+44$ 被写为 118，经原文与算术核对修正为 150。AI 草稿与纠正依据分别留在文献工作流记录中，正文使用核定后的文献主张。

认知编码先在既有真实 6 题与合成 6 题上实际调用流程，检查双维标签、连续引文和弃权。该探索采用三模型独立判断、互评与必要的仲裁；表\ref{app:pilot}保留实际有效输出和各阶段任务候选。旧规则要求三票均非空才形成多数票，并允许已执行的仲裁覆盖此前结果；其问题集、规则和模型条件与正文 t3 不同。

\begin{table}[htbp]\centering\small
\caption{早期 AI 流程探索；真实与合成各 6 题}\label{app:pilot}
\begin{tabular}{@{}lcc@{}}\toprule
阶段 & 真实 6 题 & 合成 6 题\\\midrule
独立输出有效数 & 18/18 & 18/18\\
互评输出有效数 & 17/18 & 18/18\\
仲裁输出数 & 3 & 1\\
单 Claude 候选 & 3 & 6\\
独立多数票候选 & 1 & 6\\
互评多数票候选 & 2 & 6\\
最终候选 & 2 & 6\\
互评至最终的非空改标数 & 2 & 1\\\bottomrule
\end{tabular}
\end{table}

随后小范围人工探索收集首轮独立提交 12 份，并形成第二轮 8 题、双人独立提交 16 份。分歧暴露出名词短语、清单请求、具体应用和分析等边界解释不一致。AI 据此把说明组织为主要操作、对象或条件、证据和边界四部分，提供同题推荐；第三轮两名评审分别完成 8 份推荐辅助复评。

正文表\ref{tab:human}保留独立和辅助条件下的原始统计。推荐填入、理由模板及手动填写分别为 11、4、1 份；层级与推荐一致为 13/16，标签、证据与理由逐字相同为 9/16。四段内容可识别为 15/16，形式完整性与论证充分性分别检查。最后对 4 道边界题完成用户人工核验并封存，保留原独立与辅助提交。该链条记录 AI 如何形成建议、人如何修改或采纳，以及哪些分歧仍需要保留；它不充当独立学习效果测量。

\subsection{调用记录与核查范围}
逐回合 t3 的 \path{plan.json} 保存规则、样本哈希、模型请求参数、输出上限及代码版本，\path{sample.json} 保存来源坐标。私有任务数据库保存消息与请求条件，尝试账本记录状态和可用返回；冻结清单把逐题结果与聚合分析绑定。技术失败、重试、缓存复用和逻辑任务分别保留。早期部分失败仅保存类别与用量，记录完整性按实际可得范围说明。

模型提示要求把学生文本中的指令视为研究数据，分开判断任务和贡献，证据不足时弃权，引文逐字连续匹配。程序核验字段、取值、重复键、引文位置、来源连接及统计分母；团队核查语义、文献和数学解释。模型自报概率、模型间一致性、人工推荐采纳和独立准确率具有不同含义。

AI 解题过程与采用记录的阅读入口为 \texttt{docs/AI}解题过程\texttt{.md}，其中连接论文、研究 Notebook 与工作流说明；逐轮人工与调用原始记录保存在授权档案中。论文、图表和演示保留所使用版本和来源，公开材料使用聚合统计与合成例子。学生原文、私有标识及密钥不进入公开仓库。
